% Part: computability
% Chapter: recursive-functions
% Section: pr-relations

\documentclass[../../../include/open-logic-section]{subfiles}

\begin{document}

\olfileid{cmp}{rec}{prr}
\olsection{आदिम पुनरावर्ती संबंध}


\begin{defn}
संबंध $R(\vec x)$ चे जातिबोधक फल
\[
\Char{R}(\vec x) = \left\{
  \begin{array}{ll}
  1 & \mbox{जर $R(\vec x)$ असेल} \\
  0 & \mbox{अन्यथा}
  \end{array}
\right.
\]
आदिम पुनरावर्ती असेल, तर त्या संबंधाला आदिम पुनरावर्ती म्हणतात.
\end{defn}

दुसऱ्या शब्दांत, आदिम पुनरावर्ती संबंध $R(\vec x)$ असा उल्लेख केला, तर
$\Char{R}(\vec x) = 1$ या रूपातील संबंध अभिप्रेत असतो; येथे $\Char{R}$
हे कोणत्याही आदानासाठी 1 किंवा 0 यांपैकी एक मूल्य परत करणारे आदिम
पुनरावर्ती फलन असते. उदाहरणार्थ, $x = 0$ असताना आणि केवळ तेव्हाच सत्य
असलेला संबंध $\fn{IsZero}(x)$ हा $\Char{\fn{IsZero}}$ या फलनाशी संबंधित
आहे; हे फलन आदिम पुनरावर्तन वापरून पुढीलप्रमाणे परिभाषित केले जाते:
\begin{align*}
\Char{\fn{IsZero}}(0) & = 1,\\
\Char{\fn{IsZero}}(x+1) & = 0.
\end{align*}

संबंधांचे इतर आदिम पुनरावर्ती फलनांशी संयोजन करता येते, हे स्पष्ट असावे.
म्हणून पुढील संबंधही आदिम पुनरावर्ती आहेत:
\begin{enumerate}
\item समानता संबंध, $x = y$, जो $\fn{IsZero}(\left|x -
  y\right|)$ ने परिभाषित होतो
\item लघुतर अथवा समान संबंध, $x \leq y$, जो $\fn{IsZero}(x
  \tsub y)$ ने परिभाषित होतो
\end{enumerate}


\begin{prop}
  आदिम पुनरावर्ती संबंधांचा संच बूलीय क्रियांखाली बंद आहे; म्हणजे
  $P(\vec x)$ आणि $Q(\vec x)$ आदिम पुनरावर्ती असतील, तर पुढील संबंधही
  आदिम पुनरावर्ती असतात:
  \begin{enumerate}
  \item $\lnot P(\vec x)$
  \item $P(\vec x) \land Q(\vec x)$
  \item $P(\vec x) \lor Q(\vec x)$
  \item $P(\vec x) \lif Q(\vec x)$
  \end{enumerate}
\end{prop}

\begin{proof}
  $P(\vec x)$ आणि $Q(\vec x)$ आदिम पुनरावर्ती आहेत, म्हणजे त्यांची
  जातिबोधक फले $\Char{P}$ आणि $\Char{Q}$ आदिम पुनरावर्ती आहेत, असे
  समजा. $\lnot P(\vec x)$ इत्यादींची जातिबोधक फलेही आदिम पुनरावर्ती
  आहेत हे दाखवायचे आहे.
  \[
  \Char{\lnot P}(\vec x) = \begin{cases}
    0 & \text{जर $\Char{P}(\vec x) = 1$ असेल}\\
    1 & \text{अन्यथा}
  \end{cases}
  \]
  आपण $\Char{\lnot P}(\vec x)$ ची व्याख्या $1 \tsub \Char{P}(\vec x)$
  अशी देऊ शकतो.
  \[
  \Char{P \land Q}(\vec x) = \begin{cases}
    1 & \text{जर $\Char{P}(\vec x) = \Char{Q}(\vec x) = 1$ असेल}\\
    0 & \text{अन्यथा}
  \end{cases}
  \]
  आपण $\Char{P \land Q}(\vec x)$ ची व्याख्या $\Char{P}(\vec x) \cdot
  \Char{Q}(\vec x)$ किंवा $\fn{min}(\Char{P}(\vec x), \Char{Q}(\vec
  x))$ अशी देऊ शकतो. त्याचप्रमाणे,
  \begin{align*}
    \Char{P \lor Q}(\vec x) & = \fn{max}(\Char{P}(\vec x), \Char{Q}(\vec x)) \text{ आणि}\\
    \Char{P \lif Q}(\vec x) & = \fn{max}(1 \tsub
  \Char{P}(\vec x), \Char{Q}(\vec x)).
  \end{align*}
\end{proof}

\begin{prop}
  आदिम पुनरावर्ती संबंधांचा संच परिबद्ध संख्यकीकरणाखाली बंद आहे; म्हणजे
  $R(\vec x, z)$ हा आदिम पुनरावर्ती संबंध असेल, तर पुढील संबंधही आदिम
  पुनरावर्ती असतात:
  \begin{align*}
    & \bforall{z < y}{R(\vec x, z)} \text{ आणि}\\
    & \bexists{z < y}{R(\vec x, z)}.
  \end{align*}
  प्रत्येक~$y$ पेक्षा लहान~$z$ साठी $R(\vec x, z)$ सत्य असेल तेव्हा
  आणि केवळ तेव्हाच $\bforall{z < y}{R(\vec x, z)}$ हे $\vec x$ आणि
  $y$ यांसाठी सत्य असते; $\bexists{z < y}{R(\vec x, z)}$ साठीही समरूप
  अर्थ लागू होतो.
\end{prop}

\begin{proof}
  संकेताप्रमाणे आपण $\bforall{z < 0}{R(\vec x, z)}$ हे सत्य मानतो
  (कारण~$0$ पेक्षा लहान कोणतेही $z$ नसते) आणि
  $\bexists{z < 0}{R(\vec x, z)}$ हे असत्य मानतो. परिबद्ध सार्वत्रिक
  संख्यापक सांत गुणाकार किंवा पुनरावृत्त लघुतमाप्रमाणे कार्य करतो; म्हणजे
  $P(\vec x, y) \defiff \bforall{z < y}{R(\vec x, z)}$ असेल, तर
  $\Char{P}(\vec x, y)$ ची व्याख्या अशी देता येते:
  \begin{align*}
    \Char{P}(\vec x, 0) & = 1\\
    \Char{P}(\vec x, y+1) & =
    \fn{min}(\Char{P}(\vec x, y), \Char{R}(\vec x, y)).
  \end{align*}
  परिबद्ध अस्तित्ववाची संख्यकीकरणाची व्याख्या त्याचप्रमाणे $\fn{max}$
  वापरून देता येते. पर्यायाने, सममूल्यता
  $\bexists{z < y}{R(\vec x, z)}
  \liff \lnot \bforall{z < y}{\lnot R(\vec x, z)}$ वापरून ते परिबद्ध
  सार्वत्रिक संख्यकीकरणापासून परिभाषित करता येते. उदाहरणार्थ,
  $\bexists{x \leq y}{\dots x\dots}$ या रूपातील परिबद्ध संख्यापक
  $\bexists{x < y+1}{\dots x \dots}$ याच्याशी सममूल्य आहे, हे लक्षात घ्या.
\end{proof}

\begin{prob}
  तीन-स्थानी संबंध $x \equiv y \mod n$ (भाजक~$n$ समशेषता) हा आदिम
  पुनरावर्ती आहे, हे दाखवा.
\end{prob}

आणखी एक उपयोगी आदिम पुनरावर्ती फलन म्हणजे पुढीलप्रमाणे परिभाषित केलेले
सशर्त फलन $\fn{cond}(x,y,z)$:
\begin{align*}
  \fn{cond}(x,y,z) & = \begin{cases}
  y & \text{जर $x = 0$ असेल} \\
  z & \text{अन्यथा}.
\end{cases}
\intertext{याची पुनरावर्ती व्याख्या अशी आहे:}
\fn{cond}(0,y,z) & = y,\\
\fn{cond}(x+1,y,z) & = z.
\end{align*}
आदिम पुनरावर्ती संबंधांपासून आदिम पुनरावर्ती फलनांच्या प्रकरणांनुसार
केलेल्या व्याख्यांचे समर्थन करण्यासाठी हे फलन वापरता येते:

\begin{prop}
$g_0(\vec x)$, \dots,~$g_m(\vec x)$ ही आदिम पुनरावर्ती फलने आणि $R_0(\vec
x)$, \dots, $R_{m-1}(\vec x)$ हे आदिम पुनरावर्ती संबंध असतील, तर
पुढीलप्रमाणे परिभाषित केलेले फलन $f$ देखील आदिम पुनरावर्ती असते:
\[
f(\vec x) = \begin{cases}
    g_0(\vec x) & \text{जर $R_0(\vec{x})$ असेल} \\
    g_1(\vec x) & \text{जर $R_1(\vec{x})$ असेल आणि $R_0(\vec{x})$ नसेल} \\
    \vdots & \\
    g_{m-1}(\vec x) & \text{जर $R_{m-1}(\vec{x})$ असेल आणि आधीच्या
      कोणत्याही अटी सत्य नसतील}
    \\
    g_m(\vec x) & \mbox{अन्यथा}
\end{cases}
\]
\end{prop}

\begin{proof}
  $m = 1$ असताना हे फक्त पुढीलप्रमाणे परिभाषित केलेले फलन आहे:
  \[
  f(\vec x) = \fn{cond}(\Char{\lnot R_0}(\vec x),g_0(\vec x),g_1(\vec
  x)).
  \]
  $m$ हे $1$ पेक्षा मोठे असताना या रूपातील व्याख्यांचे केवळ संयोजन करता
  येते.
\end{proof}
\end{document}
